% Generated draft: requires mathematical and source review.
\documentclass[11pt]{article}
\usepackage[margin=25mm]{geometry}
\usepackage{fontspec}
\setmainfont{Latin Modern Roman}
\usepackage{amsmath,amssymb,mathtools,unicode-math}
\setmathfont{Latin Modern Math}
\usepackage{enumitem,needspace,xurl,hyperref}
\setlength{\parindent}{0pt}
\setlength{\parskip}{5pt}
\setlength{\emergencystretch}{3em}
\newcommand{\N}{\mathbb{N}}
\newcommand{\Z}{\mathbb{Z}}
\newcommand{\Q}{\mathbb{Q}}
\newcommand{\R}{\mathbb{R}}
\newcommand{\C}{\mathbb{C}}
\newcommand{\F}{\mathbb{F}}
\newcommand{\A}{\mathbb{A}}
\newcommand{\PP}{\mathbb{P}}
\newcommand{\E}{\mathbb{E}}
\newcommand{\eps}{\varepsilon}
\newcommand{\dd}{\,\mathrm{d}}
\begin{document}
\section*{UnsolvedMath ID: 30005446}
\textbf{Source problem code:} catalogue entry
\subsection{Definitions and mathematical statement}
A percolation model specifies a graph and a probability law for declaring vertices or edges open. A percolation cluster is a connected component of the resulting open subgraph. Criticality refers to the model's critical parameter, denoted by \(p_c\); the graph, percolation model, and definition of this threshold must be specified before this notation determines a probability law. A random walk on a cluster is a stochastic process taking values in that cluster, with transition probabilities and a discrete or continuous time convention that must also be specified.

\textbf{Research question supplied by the source.} Describe the expected anomalous behaviour of random walks on percolation clusters at criticality.

Here ``anomalous behaviour'' is not assigned a mathematical definition in the supplied record. In particular, the record specifies neither a displacement observable nor a space--time scaling, limiting process, or convergence assertion. It also leaves unspecified the underlying graph and dimension, the choice or conditioning of a critical cluster, the starting point, the walk's transition rule, and whether assertions concern a fixed random environment or averages over environments. Consequently, this question does not determine a unique quantified conjecture. No particular scaling exponent, universal assertion over percolation models, or precise limiting theorem can be extracted without additional hypotheses.

The supplied background explicitly characterizes the extracted passage as contextual prose rather than a precise conjecture and states that critical clusters and anomalous random-walk scaling depend on the dimension and model. The statement above therefore preserves the research question while leaving its unresolved specifications visible.
\subsection{Short English statement}
The source asks how random walks behave anomalously on critical percolation clusters. It does not specify the model, cluster selection, walk, or meaning of anomalous scaling, so a precise conjecture requires further information. The supplied open label is a curation status, not an independently verified mathematical conclusion.
\subsection{Sources}
\begin{itemize}
\item Local registry: \texttt{data/unsolvedmath/problems.json}, permanent UnsolvedMath ID \(30005446\), source problem code \texttt{catalogue entry}, title ``Critical Random Walks on Percolation Clusters.'' The supplied record belongs to the Oberwolfach Reports Open Problems collection and gives proposed year \(2023\); no proposer is named.
\item Supplied report citation: \emph{Mini-Workshop: New Horizons in Motions in Random Media} (2023), Oberwolfach Reports, pp.\ 565--606, DOI \texttt{10.4171/owr/2023/10}. Original report link: \url{https://doi.org/10.4171/owr/2023/10}.
\item The extracted original passage associates the question with David A. Croydon's contribution, ``Random walks on random graphs in critical regimes.'' This association does not establish that Croydon proposed a particular conjecture.
\item Additional supplied literature source, labelled ``Croydon, random walks on critical percolation clusters'': \url{https://www.maths.ox.ac.uk/node/4868}.
\item Literature assessment dated \texttt{2026-08-22}: the supplied curation labels the item open and describes it as contextual prose, with dimension- and model-dependent scaling. This is supplied dated curation, not an independently verified literature assessment. The background requests independent verification of time-sensitive status and cited claims.
\item The registry labels the statement \texttt{corrected\_verified} and reports that the official report was checked directly. This is a supplied verification claim; no independent source inspection is asserted in this draft.
\end{itemize}
\end{document}
